\documentclass[11pt]{article}
\usepackage[margin=0.75in]{geometry}
\usepackage{fancyhdr}
\usepackage{array}
\usepackage{booktabs}
\usepackage{xcolor}
\usepackage{enumitem}

\pagestyle{fancy}
\fancyhf{}
\rhead{Patient: Santos, Maria E.}
\lhead{Riverside Neurology Associates}
\rfoot{Page \thepage}

\definecolor{neuroblue}{RGB}{0,70,130}

\begin{document}

\begin{center}
{\Large\bfseries\color{neuroblue} RIVERSIDE NEUROLOGY ASSOCIATES}\\[0.3em]
{\small 750 Medical Arts Building, Suite 400, Riverside, CA 92501}\\
{\small Phone: (951) 555-6200 | Fax: (951) 555-6201}\\[1em]
{\LARGE\bfseries CONSULTATION NOTE}
\end{center}

\vspace{0.5em}
\hrule height 2pt
\vspace{1em}

\section*{PATIENT INFORMATION}

\begin{tabular}{@{}p{3.5cm}p{5cm}p{3cm}p{4cm}@{}}
\textbf{Patient:} & Maria Elena Santos & \textbf{DOB:} & 03/14/1988 \\
\textbf{Date of Service:} & February 12, 2024 & \textbf{Visit Type:} & Consultation \\
\textbf{Referring MD:} & Dr. Jennifer Walsh & \textbf{Insurance:} & Hartford WC \\
\textbf{Physician:} & Sarah Goldstein, MD & & \\
\end{tabular}

\section*{REASON FOR CONSULTATION}

Evaluation of concussion following workplace fall on 02/08/2024. Referred from ER for post-concussion management.

\section*{HISTORY OF PRESENT ILLNESS}

Ms. Santos is a 35-year-old right-handed female carpenter who sustained a head injury when she fell approximately 15 feet from scaffolding on February 8, 2024. She struck her head on the concrete floor, resulting in a 3cm forehead laceration that was sutured in the ER. CT head at the time was negative for intracranial hemorrhage or skull fracture.

\textbf{Current Symptoms (4 days post-injury):}
\begin{itemize}[nosep]
\item Headache: Constant, described as pressure-like, rated 5/10
\item Dizziness: Intermittent, worse with position changes
\item Difficulty concentrating: ``Can't focus on reading or TV''
\item Fatigue: ``I sleep 12 hours and still feel tired''
\item Photophobia: Mild sensitivity to bright lights
\item Irritability: Reports feeling ``on edge''
\item Memory: Some difficulty with short-term recall
\end{itemize}

She denies loss of consciousness at the scene (per witness reports she was awake immediately after the fall, though dazed). She denies nausea/vomiting, vision changes, weakness, numbness, or seizure activity.

\section*{PAST MEDICAL HISTORY}

No prior head injuries or concussions. No history of migraines. No neurological conditions.

\section*{MEDICATIONS}

\begin{itemize}[nosep]
\item Norco 5/325mg PRN (taking 1-2x daily)
\item Ibuprofen 600mg TID
\item Flexeril 10mg QHS
\end{itemize}

\section*{PHYSICAL EXAMINATION}

\textbf{Vital Signs:} BP 128/78, HR 72, RR 14

\textbf{General:} Alert, oriented x3, appears fatigued.

\textbf{Head:} Sutured laceration to right forehead healing well. No CSF leak. No Battle's sign or raccoon eyes.

\textbf{Cranial Nerves:}
\begin{itemize}[nosep]
\item II: Visual acuity 20/20 OU, visual fields full, pupils equal and reactive
\item III, IV, VI: Extraocular movements intact, no nystagmus
\item V: Sensation intact V1-V3 bilaterally
\item VII: Face symmetric
\item VIII: Hearing grossly intact bilaterally
\item IX, X, XI, XII: Intact
\end{itemize}

\textbf{Motor:} Strength 5/5 throughout (limited testing of RUE due to wrist fracture in splint)

\textbf{Sensory:} Intact to light touch throughout

\textbf{Coordination:} Finger-to-nose intact bilaterally. Heel-to-shin intact.

\textbf{Gait:} Slightly unsteady, narrow-based. No frank ataxia.

\textbf{Cognitive Screen (MOCA):} 24/30 (Mild impairment)
\begin{itemize}[nosep]
\item Lost points: Delayed recall (2), Attention/serial 7s (2), Abstraction (1), Orientation (1)
\end{itemize}

\section*{ASSESSMENT}

\begin{enumerate}
\item \textbf{Post-Concussion Syndrome} - Patient exhibits multiple symptoms consistent with mild traumatic brain injury including headache, cognitive difficulties, fatigue, and vestibular symptoms. MOCA score of 24/30 supports mild cognitive impairment.

\item \textbf{Workplace injury} - Mechanism consistent with concussive injury.
\end{enumerate}

\section*{PLAN}

\begin{enumerate}
\item \textbf{Cognitive Rest:} Limit screen time, reading, and mentally demanding activities for the next 1-2 weeks. Gradually increase activity as tolerated.

\item \textbf{Physical Rest:} No strenuous activity. Light walking is acceptable. No return to work until cleared.

\item \textbf{Sleep Hygiene:} Maintain regular sleep schedule. It is normal to need extra sleep during recovery.

\item \textbf{Headache Management:}
\begin{itemize}[nosep]
\item Continue ibuprofen as needed but avoid overuse (rebound headache risk)
\item May add acetaminophen 1000mg as needed
\item Consider Tylenol \#3 for severe headaches
\end{itemize}

\item \textbf{Return Precautions:} Return immediately or go to ER for: severe headache, repeated vomiting, seizures, confusion, slurred speech, weakness on one side, unequal pupils.

\item \textbf{Driving:} No driving until symptoms significantly improved and cleared by this office.

\item \textbf{Work Status:} \textbf{TOTALLY DISABLED} from all work activities. Patient cannot safely perform job duties or operate machinery.

\item \textbf{Follow-up:} Return in 2 weeks for re-evaluation. If symptoms persist beyond 4 weeks, will consider neuropsychological testing and vestibular rehabilitation referral.

\item \textbf{Prognosis:} Good. Majority of concussion symptoms resolve within 2-4 weeks. Given mechanism and initial presentation, anticipate full recovery though timeline may be prolonged due to concurrent injuries.
\end{enumerate}

\vspace{2em}

\begin{tabular}{@{}p{7cm}p{7cm}@{}}
\hrulefill & \\
\textbf{Sarah Goldstein, MD} & \textbf{Date:} February 12, 2024 \\
Board Certified Neurology & \\
Riverside Neurology Associates & \\
\end{tabular}

\end{document}
